% concatenated from large-clique-8-26.tex
\documentclass[11pt,a4paper]{article}

\usepackage[numbers,sort&compress]{natbib}
\usepackage[colorlinks=true,linkcolor=blue,citecolor=magenta,urlcolor=blue]{hyperref}
\usepackage{amsmath,amssymb,amsfonts,mathrsfs,bm}
\usepackage{amsthm}
\usepackage{enumerate}
\usepackage[top=1.9cm,bottom=1.9cm,left=2.0cm,right=2.0cm]{geometry}
\usepackage{indentfirst}
\usepackage{setspace}
 %\usepackage[notref,notcite]{showkeys} % enable for draft labels

\newtheorem{theorem}{Theorem}[section]
\newtheorem{definition}[theorem]{Definition}
\newtheorem{proposition}[theorem]{Proposition}
\newtheorem{lemma}[theorem]{Lemma}
\newtheorem{remark}[theorem]{Remark}
\newtheorem{claim}[theorem]{Claim}
\newtheorem{corollary}[theorem]{Corollary}
\newtheorem{conjecture}[theorem]{Conjecture}
\newtheorem{problem}[theorem]{Problem}

\newcommand{\ex}{\operatorname{ex}}
\newcommand{\cl}{\mathcal K}
\newcommand{\Ncl}[2]{N_{#1}(#2)}
\newcommand{\dunion}{\mathbin{\dot\cup}}

\title{Large Cliques and Clique Spectral Radius in the Erd\H{o}s--S\'{o}s Problem}
\author{Xiaojun Zhao\thanks{School of Mathematics, Hunan University, Changsha 410082, P.R. China. Email: xiaojunzhao@hnu.edu.cn.}
\and Yuejian Peng\thanks{Corresponding author. School of Mathematics, Hunan University, Changsha 410082, P.R. China. Email: ypeng1@hnu.edu.cn.}}
%\date{}

\begin{document}
\setlength{\baselineskip}{14pt}
\maketitle


\begin{abstract}
For graphs \(H\) and \(F\), let \(\ex(n,H,F)\) be the maximum number of copies of \(H\) in an \(n\)-vertex \(F\)-free graph. We study this problem when \(H\) is a clique and \(F=T_t\) which is a fixed tree on \(t\) vertices. The Erd\H{o}s--S\'{o}s conjecture concerns the value of \(\ex(n,K_2, T_t)\). Gerbner and Palmer\cite{GP2026} proposed a more general conjecture: if $n=\alpha(t-1)+\beta$ and $0\le\beta\le t-2$, then the graph \(\alpha K_{t-1}\dunion K_\beta\) maximizes the number of \(r\)-cliques among all $n$-vertex $T_t$-free graphs for every \(3\le r\le t-2\).  We show that this conjecture holds for $T_t$ having at least $t-r$ leaves with a common parent, which contains the star case as a special case and recovers the sharp clique-counting result conjectured by Gan, Loh and Sudakov~\cite{gan2015maximizing} and proved by Chase~\cite{chase2020} and Chao and Dong~\cite{Chao2022}.  We also study the clique-spectral analogue. Under the same leaf-bunch condition, every $T_t$-free graph $G$ satisfies $\rho_r(G)\le\binom{t-2}{r-1}$, with equality, for $n\ge t-1$, if and only if $K_{t-1}$ is a component of $G$. Furthermore,
 we prove the conjecture for \(r=t-d\) whenever \(d\ge2\) and \(t\ge d^2-d+3\), while the case \(d=1\) is determined exactly for every \(t\). For \(d\ge2\) and \(t\ge d^2-d+3\), every \(T_t\)-free graph \(G\) satisfies \(\rho_{t-d}(G)\le\rho_{t-d}(K_{t-1})\), with equality characterized by the presence of a \(K_{t-1}\)-component.

 Our method is designed for relatively large cliques. In the leaf-poor case, after deleting edges that lie in no  $(t-d)$-clique, we study the intersection relation among \((t-d)\)-cliques and show that its equivalence classes induce the nontrivial clique-supported components; furthermore, we show each non-trivial component has at most \(t-1\) vertices. In the complementary leaf-rich case, a leaf-bunch criterion reduces the clique-counting problem to the sharp bounded-maximum-degree clique theorem.

\medskip
\noindent\emph{Keywords:} Erd\H{o}s--S\'{o}s conjecture; trees; cliques; clique spectral radius.
\end{abstract}

\section{Introduction}\label{intro}


All graphs considered in this paper are simple and undirected.
For a graph $G$, let $V(G)$ and $E(G)$ denote its vertex set and edge set, respectively.
For $S\subseteq V(G)$, let $G[S]$ denote the subgraph of $G$ induced by $S$.
For a vertex $u$ of $G$, the neighborhood of $u$ in $G$ is denoted by $N_G(u)$, or simply $N(u)$.
Let $P_t$, $K_t$, and $S_t$ denote the path, the complete graph, and the star on $t$ vertices, respectively. A copy of $K_r$ is also called an {\em $r$-clique}.
%Let $F_{t}(s,l)$ denote the graph  consisting of $l$ cliques of order $t$ sharing $s$ common vertices.
A {\em component} of $G$ is a maximal connected subgraph of $G$.

For pairwise vertex-disjoint graphs \(G_1,\ldots,G_k\), we write
$G_1\dunion\cdots\dunion G_k$ for their disjoint union. We also write \(mG\) for the disjoint union of \(m\) copies of \(G\). For example, \(\alpha K_{t-1}\dunion K_\beta\) consists of the vertex-disjoint union of \(\alpha\) copies of \(K_{t-1}\) and one copy of \(K_\beta\).


For a graph $G$ and a graph $F$,
$G$ is called {\em $F$-free} if it does not contain an isomorphic copy of $F$ as a subgraph. For an integer $n$ and graphs $H$ and $F$, let $\ex(n,H,F)$ denote the maximum number of copies of $H$ that an $n$-vertex $F$-free graph can contain; this number is called the {\em generalized Tur\'{a}n number}. An $n$-vertex $F$-free graph containing $\ex(n,H,F)$ copies of $H$ is called an {\em extremal graph}.
When $H=K_{2}$, this is the usual Tur\'{a}n number. The study of Tur\'{a}n numbers is inspired by Tur\'{a}n~\cite{t1941}. Let $T_r(n)$ denote the complete $r$-partite graph on $n$ vertices whose part sizes are as equal as possible. Tur\'{a}n's Theorem \cite{t1941} says that if $G$ is a $K_{r+1}$-free graph on $n$ vertices, then $e(G)\le e(T_r(n))$, with equality if and only if $G=T_r(n)$. A fundamental result of Erd\H{o}s, Stone and Simonovits \cite{ES46, ES66} gives that $\ex(n,K_2,F)=e(T_r(n))+o(n^2)$, where $\chi(F)=r+1\ge3$. This determines the asymptotic value of $\ex(n,K_2,F)$ for every non-bipartite graph $F$. Thus a remaining challenge is to determine Tur\'{a}n numbers of bipartite graphs.

Determining the Tur\'{a}n number of a tree is an important part of this theory. Clearly, every graph with $n$ vertices and more than $\frac{n(t-2)}{2}$ edges contains a star with $t$ vertices.
Erd\H{o}s and S\'{o}s \cite{erdos1962} in 1962 conjectured that every graph with $n$ vertices and more than $\frac{n(t-2)}{2}$ edges contains every tree with $t$ vertices as a subgraph.
The conjectured bound is tight for every $t\in\mathbb N$, as seen by considering a graph consisting of a vertex-disjoint union of copies of $K_{t-1}$ if $t-1$ divides $n$.
 The Erd\H{o}s--Gallai theorem \cite{1959Erdos} shows that the conjecture is true for a path.
 Ajtai, Koml\'{o}s and Szemer\'{e}di \cite{ajtai1994} announced a proof of the Erd\H{o}s--S\'{o}s conjecture for large $t$, and notable works (for example, see \cite{Brandt1996,goerlich2016,rozhon2019,besomi2019,Fan2013,Fan2007,Fan2019,sacle1997}) have verified the Erd\H{o}s--S\'{o}s conjecture in various special cases, but the conjecture remains widely open.








%In this paper, we focus on  $ex(n,H,\mathcal{F})$ for $H\neq K_{2}$.
The study of generalized Tur\'{a}n numbers
 was initiated by Zykov \cite{zykov1949}, who determined $\ex(n,K_{r},K_{t})$ exactly, and
 a systematic study of the generalized Tur\'{a}n problem $\ex(n,H,\mathcal{F})$ was given by Alon and Shikhelman in \cite{alon2016}. For a recent survey of generalized Tur\'{a}n problems, see Gerbner and Palmer \cite{GP2026}.
 %In this paper, we focus on $\ex(n,K_r,T)$ for a tree $T$.




For the path,
Luo \cite{luo2018maximum} generalized the Erd\H{o}s--Gallai theorem and proved $\ex(n,K_{r},P_{t})\leq \frac{n}{t-1}\binom{t-1}{r}$.
Chakraborti and Chen \cite{chakraborti2024exact} strengthened Luo's result by obtaining the exact value for every $n$ and characterizing all extremal graphs.

 For the star, Cutler and Radcliffe \cite{cutler2014maximum} determined the maximum number of complete subgraphs in $S_{t}$-free graphs $G$ on $n$ vertices and characterized all extremal graphs.
A related question is to determine $\ex(n,K_r,S_t)$.
Several works provide values for small $n$; see \cite{engbers2014counting,alexander2012,alexander2016,Law2013}.
Gan, Loh and Sudakov \cite{gan2015maximizing} conjectured that, if $n=\alpha(t-1)+\beta$ with $0\leq\beta\leq t-2$, then
\[
\ex(n,K_r,S_t)= \alpha\binom{t-1}{r}+\binom{\beta}{r}.
\]
Chase \cite{chase2020} confirmed the conjecture by showing that it holds for $r=3$ and combining the result of Gan, Loh and Sudakov \cite{gan2015maximizing} that the conjecture holds for $r+1$ if it holds for $r$.  Chao and Dong \cite{Chao2022} later gave a unified proof for all clique orders $r\ge3$.



 Paths and stars have special structures. It is therefore natural to ask the same question for a general tree. In 2026, Gerbner and Palmer \cite{GP2026} proposed the following general conjecture for $r$-cliques with $3\leq r\leq t-2$. Independently, we had formulated the same conjecture in an unpublished first draft of this paper in May 2025.



\begin{conjecture}\label{C1}[Gerbner and Palmer \cite{GP2026}]
Let $\alpha\geq 1, 0\leq \beta\leq t-2, 3\leq r\leq t-2$ be integers and $n=\alpha (t-1)+\beta$.
    For each $t$-vertex tree $T_{t}$, we have
\[
\ex(n,K_r,T_t)=\alpha\binom{t-1}{r}+\binom{\beta}{r}.
\]
\end{conjecture}

The graph \(\alpha K_{t-1}\dunion K_\beta\) is \(T_t\)-free and therefore gives $\ex(n,K_r,T_t)\ge \alpha\binom{t-1}{r}+\binom{\beta}{r}$.


When we were working on the problem in \cite{ZP1} of determining $\ex(n,K_3,\{C_5,C_{10},C_{15},C_{20},\ldots\})$, we realized that the method designed for counting the maximum number of copies of $K_3$ in a graph forbidding cycles whose lengths are divisible by $5$ might be modified to determine $\ex(n,K_r,T_t)$ for large $r$ and a $t$-vertex tree $T_t$. In the first draft of this paper, written in May 2025, we formulated Conjecture~\ref{C1} independently and proved the case $r=t-2$. We subsequently extended our argument to the cases $r=t-3$ and $r=t-4$. These drafts were not posted online, since we continued working toward more general results using our method. Very recently, Zhou and Yuan~\cite{ZhouYuan2026} posted an independent proof of Conjecture~\ref{C1} for $r=t-2$ and for $r=t-3\ge5$, together with a characterization of all extremal graphs.  In the current version of this paper, we determine $\ex(n,K_{t-d},T_t)$ for $d=O(\sqrt{t})$ and characterize all extremal graphs (Theorem~\ref{T5}). We also show that conjecture holds for an $t$-vertex tree having at least $t-r$ leaves with a common parent.

In addition to the known results on paths and stars, Gerbner, Methuku, and Palmer \cite{GMP2020} proved that if the Erd\H{o}s--S\'{o}s conjecture holds for \(T_t\) and all its subtrees, then
\[
\ex(n,K_r,T_t)\le \frac{n}{t-1}\binom{t-1}{r}
\]
for \(3\le r\le t-3\). This bound is also sharp when \((t-1)\mid n\). Gerbner \cite{gerbner2023doublestars} also showed that the conjecture holds for double stars. Letzter \cite{letzter2019} determined the order of magnitude of \(\ex(n,H,T)\) for every fixed graph \(H\) and tree \(T\).


 We show that this conjecture holds for $T_t$ having at least $t-r$ leaves with a common parent, which contains the star case as a special case.   We also study the clique-spectral analogue. Under the same leaf-bunch condition, every $T_t$-free graph $G$ satisfies $\rho_r(G)\le\binom{t-2}{r-1}$, with equality, for $n\ge t-1$, if and only if $K_{t-1}$ is a component of $G$. Furthermore,
 we prove the conjecture for \(r=t-d\) whenever \(d\ge2\) and \(t\ge d^2-d+3\), while the case \(d=1\) is determined exactly for every \(t\). We show that for \(d\ge2\) and \(t\ge d^2-d+3\), every \(T_t\)-free graph \(G\) satisfies \(\rho_{t-d}(G)\le\rho_{t-d}(K_{t-1})\), with equality characterized by the presence of a \(K_{t-1}\)-component. Let us state our results precisely in the following two subsections.









\subsection{Clique-counting results}\label{counting-results}

The case $r=t-1$ is straightforward, and we include the proof for completeness. We first observe the following.

\begin{lemma}\label{isolation}
If a \(T_t\)-free graph contains a copy \(Q\) of \(K_{t-1}\), then \(Q\) is a component.
\end{lemma}


\begin{proof}
Let \(\ell\) be a leaf of \(T_t\) and let \(p\) be its neighbor. If \(xy\in E(G)\) with \(x\in V(Q)\) and \(y\notin V(Q)\), embed \(T_t-\ell\) into \(Q\) with \(p\) mapped to \(x\), and map \(\ell\) to \(y\). This gives a copy of \(T_t\), a contradiction.
\end{proof}

\begin{theorem}\label{T1}
Let \(T_t\) be a \(t\)-vertex tree. Then
\[
\ex(n,K_{t-1},T_t)=\left\lfloor\frac{n}{t-1}\right\rfloor.
\]
The extremal graphs are exactly
\[
\left\lfloor\frac{n}{t-1}\right\rfloor K_{t-1}\dunion H,
\]
where \(H\) is any graph on \(n\bmod(t-1)\) vertices.
\end{theorem}

\begin{proof}


 By Lemma \ref{isolation}, the number of \((t-1)\)-cliques is at most \(\lfloor n/(t-1)\rfloor\). Thus, an extremal graph consists of $\lfloor n/(t-1)\rfloor$ copies of $K_{t-1}$ as components, together with an arbitrary graph on the remaining $n\bmod(t-1)$ vertices.
\end{proof}



\begin{definition}\label{leaf-bunch-parameter}
For a tree $T$, define
\[
\lambda(T):=\max_{v\in V(T)}|\{u\in N_T(v):d_T(u)=1\}|.
\]
Thus $\lambda(T)$ is the maximum number of leaves having a common parent.
\end{definition}


The following theorem isolates the clique-counting consequence of a large leaf bunch. Its clique-spectral counterpart is stated in Subsection~\ref{spectral-results}; the two theorems are proved together in Section~\ref{main-proof}.

\begin{theorem}\label{leaf-bunch-counting-theorem}
Let $T_t$ be a $t$-vertex tree and $3\le r\le t-1$. Suppose that $T_t$ has at least $t-r$ leaves with a common parent.
Let
\[
n=\alpha(t-1)+\beta,
\qquad \alpha\ge0,
\qquad 0\le\beta\le t-2.
\]
Then
\begin{equation}\label{leaf-bunch-counting}
\ex(n,K_r,T_t)
=
\alpha\binom{t-1}{r}+\binom{\beta}{r}.
\end{equation}
The extremal graphs are exactly
\[
G=\alpha K_{t-1}\dunion H,
\]
where $H=K_\beta$ if $\beta\ge r$, while $H$ is an arbitrary graph on $\beta$ vertices if $\beta<r$.
\end{theorem}

Theorem~\ref{leaf-bunch-counting-theorem} gives a direct family of trees for which Conjecture~\ref{C1} holds. Indeed, for  $3\le r\le t-2$, if $T_t$ has at least $t-r$ leaves with a common parent, then $\lambda(T_t)\ge t-r$, and Theorem~\ref{leaf-bunch-counting-theorem} gives exactly the generalized Tur\'{a}n number predicted by Conjecture~\ref{C1}, together with the extremal graphs.  In particular, Theorem~\ref{leaf-bunch-counting-theorem} contains the star case as a special case and recovers the sharp clique-counting result conjectured by Gan, Loh and Sudakov~\cite{gan2015maximizing} and proved by Chase~\cite{chase2020} and Chao and Dong~\cite{Chao2022}. In this sense, the leaf-bunch theorem extends the bounded-maximum-degree phenomenon underlying the star case to a broader class of trees.

We now state our main counting theorem.

\begin{theorem}\label{T5}
Let \(d\ge 2\) and \(t\ge d^2-d+3\), and let \(T_t\) be a \(t\)-vertex tree. Let
\[
n=\alpha(t-1)+\beta,
\qquad \alpha\ge1,
\qquad 0\le\beta\le t-2.
\]
Then
\[
\ex(n,K_{t-d},T_t)
=\alpha\binom{t-1}{t-d}+\binom{\beta}{t-d}.
\]
The extremal graphs are exactly the following:
\begin{enumerate}[(i)]
\item if \(0\le\beta<t-d\), the graph is \(\alpha K_{t-1}\dunion H\), where \(H\) is any graph on \(\beta\) vertices;
\item if \(t-d\le\beta\le t-2\), the graph is \(\alpha K_{t-1}\dunion K_\beta\).
\end{enumerate}
\end{theorem}

The proof will be given in Section 2.2. It  has two structural steps. First, a tree-separation argument shows that two \((t-d)\)-cliques with a medium-sized intersection already contain every tree on \(t\) vertices. Thus, in a \(T_t\)-free graph, intersecting large cliques have a large overlap and clique intersection becomes transitive. After deleting edges that lie in no counted clique, the intersection classes are precisely the nontrivial components. Second, Lemma~\ref{component-bound} gives a uniform bound of \(t-1\) on the orders of these components. The counting theorem then follows from a discrete convexity argument.

\subsection{Clique-spectral results}\label{spectral-results}

We next introduce the clique-spectral setting and state all spectral results used later. Let $G$ be a graph with $V(G)=\{v_1,\ldots,v_n\}$. The \emph{adjacency matrix} of $G$ is the $n\times n$ matrix
\[
A(G)=(a_{ij}),
\qquad
a_{ij}=\begin{cases}
1,& v_iv_j\in E(G),\\
0,& v_iv_j\notin E(G).
\end{cases}
\]
Since $G$ is simple and undirected, $A(G)$ is a symmetric $(0,1)$-matrix with zero diagonal. Its largest eigenvalue is called the \emph{adjacency spectral radius} of $G$ and is denoted by $\rho(G)$.   For $n>k$, let
\[
S_{n,k}=K_k\vee \overline{K}_{n-k},
\]
and let $S_{n,k}^{+}$ be the graph obtained from $S_{n,k}$ by adding one edge inside the independent set. Nikiforov~\cite{Nikiforov2010paths} conjectured that these two graphs give the sharp adjacency-spectral thresholds for forcing all trees of the corresponding orders. Cioab\u{a}, Desai and Tait~\cite{CioabaDesaiTait2023} confirmed this conjecture.

\begin{theorem}[Cioab\u{a}--Desai--Tait~\cite{CioabaDesaiTait2023}]\label{CDT-spectral-Erdos-Sos}
Let $k\ge2$ be fixed. For all sufficiently large $n$, the following statements hold.
\begin{enumerate}[(i)]
\item If an $n$-vertex graph $G$ satisfies
\[
\rho(G)\ge \rho(S_{n,k}),
\]
then either $G$ contains every tree on $2k+2$ vertices or $G\cong S_{n,k}$.
\item If an $n$-vertex graph $G$ satisfies
\[
\rho(G)\ge \rho(S_{n,k}^{+}),
\]
then either $G$ contains every tree on $2k+3$ vertices or $G\cong S_{n,k}^{+}$.
\end{enumerate}
\end{theorem}

Theorem~\ref{CDT-spectral-Erdos-Sos} corresponds to the case $r=2$ of the $r$-clique spectral radius defined below. For an integer $r\ge2$, let $\mathcal C_r(G)$ denote the family of $r$-cliques of a graph $G$.
Following Liu and Bu~\cite{LiuBu2023}, the \emph{$r$-clique tensor} of an $n$-vertex graph $G$ is the order-$r$, dimension-$n$ symmetric tensor
\[
\mathcal A_r(G)=(a_{i_1\ldots i_r}),
\]
where
\[
a_{i_1\ldots i_r}=
\begin{cases}
\dfrac{1}{(r-1)!},& \{i_1,\ldots,i_r\}\in\mathcal C_r(G),\\[2mm]
0,&\text{otherwise}.
\end{cases}
\]
The spectral radius of $\mathcal A_r(G)$ is called the \emph{$r$-clique spectral radius} of $G$ and is denoted by $\rho_r(G)$  which can be viewed as the spectral  counterpart of counting $r$-cliques.
By the variational characterization of the spectral radius of a nonnegative symmetric tensor~\cite{Qi2013},
\begin{equation}\label{spectral-variational}
\rho_r(G)=
\max\left\{
 r\sum_{\{i_1,\ldots,i_r\}\in\mathcal C_r(G)}x_{i_1}\cdots x_{i_r}:
 x_i\ge0,\ \sum_{i=1}^n x_i^r=1
\right\}.
\end{equation}
For a vertex $v$, write
\[
c_r(v)=|\{K\in\mathcal C_r(G):v\in V(K)\}|.
\]
For every admissible vector in~\eqref{spectral-variational}, the arithmetic--geometric mean inequality gives
\[
r\prod_{v\in K}x_v\le \sum_{v\in K}x_v^r.
\]
Summing over all $r$-cliques, we obtain the useful local bound
\begin{equation}\label{local-clique-degree}
\rho_r(G)\le \max_{v\in V(G)}c_r(v).
\end{equation}
The $2$-clique tensor is the adjacency matrix, and hence $\rho_2(G)=\rho(G)$.

The spectral counterpart of Theorem~\ref{T1} is also immediate.

\begin{theorem}\label{spectral-T1}
Let $T_t$ be a $t$-vertex tree and let $G$ be a $T_t$-free graph. Then
\[
\rho_{t-1}(G)\le \rho_{t-1}(K_{t-1})=1.
\]
If $|V(G)|\ge t-1$, equality holds if and only if $K_{t-1}$ is a component of $G$.
\end{theorem}

\begin{proof}
By Lemma \ref{isolation},  distinct copies of $K_{t-1}$ in $G$ are vertex-disjoint. For every vector $x\ge0$ with $\sum_v x_v^{t-1}=1$, the arithmetic--geometric mean inequality gives, for each $(t-1)$-clique $Q$,
\[
(t-1)\prod_{v\in V(Q)}x_v\le \sum_{v\in V(Q)}x_v^{t-1}.
\]
Since the $(t-1)$-cliques are vertex-disjoint, summing these inequalities and using~\eqref{spectral-variational} gives $\rho_{t-1}(G)\le1$. If $G$ contains a copy of $K_{t-1}$, equality is attained by the uniform vector supported on that clique. Hence equality is equivalent to the presence of a $K_{t-1}$-component.
\end{proof}

Recall the parameter $\lambda(T)$ from Definition~\ref{leaf-bunch-parameter}. The clique-counting form of the leaf-bunch argument was stated in Theorem~\ref{leaf-bunch-counting-theorem}. We now state its clique-spectral counterpart.

\begin{theorem}\label{leaf-bunch-spectral-theorem}
Let $T_t$ be a $t$-vertex tree and $3\le r\le t-1$. Suppose that $T_t$ contains at least $t-r$ vertices sharing a common parent.
Let $G$ be an $n$-vertex $T_t$-free graph. Then
\begin{equation}\label{leaf-bunch-criterion-bound}
\rho_r(G)\le\binom{t-2}{r-1}.
\end{equation}
If $n\ge t-1$, equality holds if and only if $K_{t-1}$ is a component of $G$.
\end{theorem}

We now state the clique-spectral result in the same $t-d$ form as Theorem~\ref{T5}.

\begin{theorem}\label{spectral-T5}
Let $d\ge2$ and $t\ge d^2-d+3$, and let $T_t$ be a $t$-vertex tree. Then every $T_t$-free graph $G$ satisfies
\[
\rho_{t-d}(G)
\le \rho_{t-d}(K_{t-1})
=\binom{t-2}{t-d-1}
=\binom{t-2}{d-1}.
\]
If $|V(G)|\ge t-1$, equality holds if and only if
\[
G=K_{t-1}\dunion H,
\]
where $H$ is an arbitrary $T_t$-free graph.
\end{theorem}

The classical Erd\H{o}s--S\'{o}s conjecture has the standard tightness construction given, when $t-1$ divides $n$, by the disjoint union of copies of $K_{t-1}$. By contrast, the exceptional graphs $S_{n,k}$ and $S_{n,k}^{+}$ in the adjacency-spectral theorem are connected join-type graphs with a small complete core and a large independent part. Nikiforov~\cite{Nikiforov2010paths} proved that $S_{n,k}$ and $S_{n,k}^{+}$ are the adjacency-spectral extremal graphs for forbidding $P_{2k+2}$ and $P_{2k+3}$, respectively, when $n$ is sufficiently large. Thus the adjacency-spectral formulation does not simply replace the number of edges in the classical tightness construction by its spectral radius.

The relation between the counting and spectral problems is more direct in our large-clique setting. Theorem~\ref{T5} shows that the $K_{t-d}$-counting extremal graph consists of as many complete blocks of order $t-1$ as possible, together with one possible smaller complete block. Theorem~\ref{spectral-T5} has the same local extremal block $K_{t-1}$, but only one such block is needed: clique counts add over components, whereas clique-spectral radius takes the maximum over clique-supported components.

Section~\ref{main-proof} contains all proofs not already given here. Section~\ref{preliminaries} proves the  structural lemmas. Section~\ref{counting-proof} gives the discrete optimization lemma, proves Theorems~\ref{leaf-bunch-counting-theorem} and~\ref{leaf-bunch-spectral-theorem} together, and then proves Theorem~\ref{T5}. Section~\ref{spectral-proofs} proves the remaining spectral lemmas and proves Theorem~\ref{spectral-T5}. Section~\ref{concluding-remarks} contains concluding remarks and open problems.

Conjecture~\ref{C1} is a clique-counting version of the Erd\H{o}s--S\'{o}s conjecture.
Our approach is not an edge-counting argument. It is a separate method for large cliques. We delete every edge not contained in an $r$-clique, study how $r$-cliques intersect in an extremal graph, and define a relation on the set of all $r$-cliques based on whether they intersect. Based on the intersection property of $r$-cliques, we show that the intersection relation is an equivalence relation and that the subgraph induced by each equivalence class is a component. Thus we get structural properties of an extremal graph.



\section{Proofs of the main results}\label{main-proof}

\subsection{Structural lemmas}\label{preliminaries}

We prove the structural lemmas needed in both the counting and clique-spectral arguments. A nontrivial component is a component with at least two vertices.





We will use the following tree decomposition due to Brennan~\cite{Brennan2016,Brennan2017}.

\begin{definition}\label{tree-division-definition}
Let \(T\) be a tree of order \(n\). A \emph{\((1/3,2/3)\)-division} of \(T\) consists of a vertex \(z\in V(T)\) and two vertex-disjoint forests \(K,H\subseteq T-z\) such that
\[
V(K)\dunion V(H)=V(T)\setminus\{z\},
\]
there is no edge between \(V(K)\) and \(V(H)\), and
\[
\frac{n-1}{3}\le |V(K)|\le\frac{n-1}{2},\qquad
\frac{n-1}{2}\le |V(H)|\le\frac{2(n-1)}{3}.
\]
\end{definition}

\begin{lemma}[Brennan~\cite{Brennan2016,Brennan2017}]\label{tree-division}
Every tree has a \((1/3,2/3)\)-division.
\end{lemma}

Lemma \ref{tree-division} has the following  direct implication.
\begin{corollary}[Brennan~\cite{Brennan2016,Brennan2017}]\label{separation}
Let \(d\ge2\), and let \(T\) be a tree on at least \(3d-1\) vertices. Then there are disjoint sets \(X,Y\subseteq V(T)\), each of order \(d\), with no edge between \(X\) and \(Y\).
\end{corollary}


\begin{lemma}\label{two-clique}
Let \(d\ge2\), \(t\ge3d-1\), and \(1\le s\le t-2d\). If \(A\) and \(B\) are cliques of order \(t-d\) with \(|V(A)\cap V(B)|=s\), then \(G[V(A)\cup V(B)]\) contains every tree on \(t\) vertices.
\end{lemma}

\begin{proof}
We use induction on \(t\), with \(d\) and \(s\) fixed. Put
\[
t_0=\max\{3d-1,2d+s\}.
\]
First consider \(t=t_0\).

If \(s\le d-1\), then \(t_0=3d-1\). By Lemma~\ref{tree-division}, let \(z\) and the two forests \(K,H\) form a \((1/3,2/3)\)-division of a $t$-vertex tree \(T\). Put \(a=|V(K)|\) and \(b=|V(H)|\). Since \(t=3d-1\), the bounds in Definition~\ref{tree-division-definition} give \(d\le a,b\le2d-2\) and \(a+b=3d-2\). The private part of either clique has order
\[
c=t_0-d-s=2d-1-s=2d-2-(s-1).
\]
Move \(s-1\) additional vertices from the two forests into a separator set \(S\) containing \(z\), choosing at least \(\max\{0,a-c\}\le s-1\) vertices from \(K\) and at least \(\max\{0,b-c\}\le s-1\) from \(H\). This is possible: if both $a-c$ and $b-c$ are positive, their sum is \(2s-d\le s-1\), and if only one is positive, it is at most \(s-1\). Every component of \(T-S\) remains in one of the two forests, and the number of vertices assigned to either side is at most \(c\). Map \(S\) into \(V(A)\cap V(B)\), and map the remaining vertices of \(K\) and \(H\) into \(V(A)\setminus V(B)\) and \(V(B)\setminus V(A)\), respectively, which embeds $T$ into \(G[V(A)\cup V(B)]\).

If \(s\ge d\), then \(t_0=2d+s\). By Corollary \ref{separation}, there are disjoint \(d\)-sets \(X,Y\subseteq V(T)\) with no edge between them. Map \(X\) and \(Y\) into the two private parts, each of order \(d\), and map the remaining \(s\) vertices into \(V(A)\cap V(B)\). Since there is no edge between \(X\) and \(Y\), this is an embedding of $T$ into \(G[V(A)\cup V(B)]\).

Now let \(t>t_0\), and let \(p\) be a leaf of \(T\). Choose private vertices \(a\in V(A)\setminus V(B)\) and \(b\in V(B)\setminus V(A)\), and put \(A_0=A-a\) and \(B_0=B-b\). These are cliques of order \((t-1)-d\) with intersection size \(s\), and \(t-1\ge t_0\). By induction, \(T-p\) can be embedded into \(A_0\cup B_0\). If the neighbor of \(p\) is mapped into \(A_0\), map \(p\) to \(a\); otherwise map it to \(b\). This completes the induction.
\end{proof}


We now use Lemma~\ref{two-clique} to describe how large cliques can intersect.

\begin{lemma}\label{large-overlap}
Let \(d\ge2\) and \(t\ge 3d-1\). Let \(T_t\) be a \(t\)-vertex tree, and let \(G\) be \(T_t\)-free. If \(K\) and \(K'\) are distinct intersecting \((t-d)\)-cliques in \(G\), then
\[
t-2d+1\le |V(K)\cap V(K')|\le t-d-1.
\]
\end{lemma}

\begin{proof}
Put \(s=|V(K)\cap V(K')|\). If \(1\le s\le t-2d\), then Lemma~\ref{two-clique} applies and shows that \(G[V(K)\cup V(K')]\) contains every \(t\)-vertex tree, a contradiction. Hence \(s\ge t-2d+1\). The upper bound follows from \(K\ne K'\).
\end{proof}

\begin{lemma}\label{large-transitivity}
Let \(\cl\) be the family of \((t-d)\)-cliques in a \(T_t\)-free graph, where \(d\ge2\) and \(t\ge3d-1\). Then the relation
\[
R\sim R'
\quad\Longleftrightarrow\quad
V(R)\cap V(R')\ne\emptyset
\]
is an equivalence relation on \(\cl\).
\end{lemma}

\begin{proof}
Only transitivity needs proof. Suppose, for a contradiction, that \(R_1\sim R_2\), \(R_2\sim R_3\), and \(V(R_1)\cap V(R_3)=\emptyset\). By Lemma~\ref{large-overlap},
\[
\begin{aligned}
t-d=|V(R_2)|
&\ge |V(R_1)\cap V(R_2)|+|V(R_2)\cap V(R_3)|\\
&\ge2(t-2d+1).
\end{aligned}
\]
Thus \(t\le3d-2\), contradicting \(t\ge3d-1\).
\end{proof}








\begin{lemma}\label{pg}
Let \(G\) be a graph in which every edge is contained in an \(r\)-clique. Let \(\cl\) be the set of all \(r\)-cliques in \(G\), and define a relation on \(\cl\) by \(R\sim R'\) if and only if \(V(R)\cap V(R')\ne\emptyset\). If \(\sim\) is an equivalence relation with classes \(\cl_1,\ldots,\cl_s\), then
\[
G[V(\cl_1)],\ldots,G[V(\cl_s)]
\]
are precisely the nontrivial components of \(G\). Here \(V(\cl_i)\) is the union of the vertex sets of the cliques in \(\cl_i\).
\end{lemma}

\begin{proof}
Within one equivalence class, any two cliques intersect, so \(G[V(\cl_i)]\) is connected. Distinct classes have disjoint vertex sets. If an edge \(xy\) joined \(V(\cl_i)\) to \(V(\cl_j)\), then \(xy\) would lie in an \(r\)-clique \(R\). The clique \(R\) would intersect a clique in \(\cl_i\) and a clique in \(\cl_j\), forcing \(i=j\). Thus no such edge exists.
\end{proof}

\begin{lemma}\label{component-bound}
Let $d\ge2$, let $T_t$ be a $t$-vertex tree, and put $r=t-d$. For a $T_t$-free graph $G$, let $G_r$ be the spanning subgraph of $G$ whose edges are precisely those contained in at least one copy of $K_r$. If
\[
\lambda(T_t)\le d-1
\qquad\text{and}\qquad
 t\ge d^2-d+3,
\]
then every nontrivial component of $G_r$ has order at most $t-1$.
\end{lemma}

\begin{proof}
Since
\[
(d^2-d+3)-(3d-1)=(d-2)^2\ge0,
\]
we have $t\ge3d-1$. By Lemma~\ref{large-transitivity}, intersection is therefore an equivalence relation on the family of $r$-cliques of $G_r$, and Lemma~\ref{pg} identifies the corresponding equivalence classes with the nontrivial components of $G_r$.

Suppose, for a contradiction, that a nontrivial component $C$ has at least $t$ vertices. Fix an $r$-clique $R$ in $C$ and choose distinct vertices
\[
u_1,\ldots,u_d\in V(C)\setminus V(R).
\]
Every vertex of $C$ lies in an $r$-clique from the equivalence class corresponding to $C$. Hence, for each $i\in[d]$, there is an $r$-clique $R_i$ containing $u_i$ and intersecting $R$. By Lemma~\ref{large-overlap},
\begin{equation}\label{component-Ai}
A_i:=N_R(u_i),
\qquad
|A_i|\ge |V(R_i)\cap V(R)|\ge t-2d+1=r-d+1.
\end{equation}
We shall embed $T_t$ into $G[V(R)\cup\{u_1,\ldots,u_d\}]$.

Let $\ell$ be the number of leaves of $T_t$.

\smallskip
\noindent\emph{Case 1: $\ell\ge d$.}

Choose any $d$ leaves, denote them by $v_1,\ldots,v_d$, map $v_i$ to $u_i$, and let $P$ be the set of their distinct parents. Put $p=|P|$. Since $\lambda(T_t)\le d-1$, we have $p\ge2$. For $y\in P$, define
\[
I(y)=\{i\in[d]:y\text{ is the parent of }v_i\},
\qquad
B_y=\bigcap_{i\in I(y)}A_i.
\]
Since the positive integers $|I(y)|$, $y\in P$, sum to $d$, we have $|I(y)|\le d-p+1$. By~\eqref{component-Ai},
\[
|B_y|\ge r-|I(y)|(d-1)
\ge t-d-(d-p+1)(d-1).
\]
Moreover,
\[
\begin{aligned}
|B_y|-p
&\ge t-d^2-d+1+p(d-2).
\end{aligned}
\]
Since $p\ge2$ and $t\ge d^2-d+3$, the right-hand side is at least
\[
t-d^2+d-3\ge0.
\]
Hence $|B_y|\ge p$ for every $y\in P$. Thus,  there are distinct vertices $b_y\in B_y$. Map each parent $y$ to $b_y$ and extend the map to a bijection from $V(T_t)\setminus\{v_1,\ldots,v_d\}$ onto $V(R)$. All tree edges are preserved, yielding a copy of $T_t$, a contradiction.

\smallskip
\noindent\emph{Case 2: $\ell<d$.}

Since $l\ge 2$, we have $d\ge3$. Let $b$ be the number of vertices of degree at least $3$ and let $h$ be the number of vertices whose degree is not $2$. The identity
\[
\ell=2+\sum_{v:\,d_{T_t}(v)\ge3}\bigl(d_{T_t}(v)-2\bigr)
\]
implies $b\le\ell-2$, and hence
\[
h=\ell+b\le2\ell-2.
\]
Suppress all degree-$2$ vertices. The resulting skeleton is a tree on $h$ vertices, whose $h-1$ edges correspond to the maximal bare paths of $T_t$.

On every maximal bare path, mark the degree-$2$ vertex closest to each end whenever it exists; if an end is a leaf, also mark the second degree-$2$ vertex from that end whenever it exists. Thus at most $2(h-1)+\ell$ degree-$2$ vertices are marked. If $M$ denotes the number of unmarked degree-$2$ vertices, then
\[
M\ge t-h-2(h-1)-\ell=t-3h+2-\ell\ge t-7\ell+8.
\]
From each remaining path choose every third unmarked vertex. The marking at the ends guarantees that choices from distinct bare paths are also at distance at least $3$, and no chosen vertex is within distance $2$ of a leaf. Moreover,
\[
\begin{aligned}
M-[3(d-\ell)-2]
&\ge t-3d-4\ell+10\\
&\ge(d^2-d+3)-3d-4(d-1)+10\\
&=d^2-8d+17=(d-4)^2+1>0.
\end{aligned}
\]
Thus we can choose $d-\ell$ degree-$2$ vertices $x_1,\ldots,x_{d-\ell}$ which are pairwise at distance at least $3$, and whose neighbors are all distinct and are different from all parents of leaves.

Map all $\ell$ leaves and the vertices $x_1,\ldots,x_{d-\ell}$ bijectively to $u_1,\ldots,u_d$. Re-indexing if necessary, suppose the leaves are mapped to $u_1,\ldots,u_\ell$. Let $P$ be the set of distinct parents of the leaves and put $p=|P|$. For $y\in P$, define
\[
I(y)=\{i\in[\ell]:y\text{ is the parent of the leaf mapped to }u_i\},
\qquad B_y=\bigcap_{i\in I(y)}A_i.
\]
Since $|I(y)|\le\ell-p+1$ and $\ell\le d-1$,
\[
\begin{aligned}
|B_y|-p
&\ge t-d-(\ell-p+1)(d-1)-p\\
&\ge t-d-(d-p)(d-1)-p\\
&\ge p(d-2)-d+3\ge0.
\end{aligned}
\]
For each selected degree-$2$ vertex mapped to $u_j$, prescribe the candidate set $A_j$ for each of its two neighbors. The total number of prescribed neighbor-images is
\[
p+2(d-\ell)\le2d-\ell\le2d-2.
\]
On the other hand, by~\eqref{component-Ai},
\[
|A_j|\ge t-2d+1\ge2d-2,
\]
because
\[
t-2d+1-(2d-2)\ge d^2-5d+6=(d-2)(d-3)\ge0.
\]
%Hall's condition follows. Indeed, a subfamily containing one of the sets $A_j$ has union of size at least $2d-2$, which is at least the total number of prescribed neighbor-images; a subfamily consisting only of the sets $B_y$ has union of size at least $p$, and there are only $p$ such sets.
Therefore all parents of leaves and all neighbors of the selected degree-$2$ vertices can be mapped to distinct vertices of $R$ in their prescribed candidate sets. Extend this map to a bijection from the remaining $t-d$ vertices of $T_t$ onto $V(R)$. Again all tree edges are preserved, giving a copy of $T_t$, a contradiction.

Both cases are impossible, so every nontrivial component of $G_r$ has order at most $t-1$.
\end{proof}



\subsection{Proofs of Theorems~\ref{leaf-bunch-counting-theorem}, \ref{leaf-bunch-spectral-theorem} and  \ref{T5}}\label{counting-proof}

We first prove the discrete optimization lemma used in the counting argument.

\begin{lemma}\label{opt}
Let \(M>r\ge2\), let \(n=\alpha M+\beta\) with \(\alpha\ge1\) and \(0\le\beta<M\), and let \(n_1,\ldots,n_k\) be positive integers such that \(n_i\le M\), each \(n_i\) is either \(1\) or at least \(r\), and \(\sum_i n_i\le n\). Then
\[
\sum_{i=1}^k\binom{n_i}{r}\le \alpha\binom Mr+\binom\beta r.
\]
Equality holds only when the parts consist of \(\alpha\) copies of \(M\), together with one part of size \(\beta\) if \(\beta\ge r\), or \(\beta\) copies of $1$ if \(\beta<r\).
\end{lemma}

\begin{proof}
Set \(f(k)=\binom{k}{r}\), with \(f(k)=0\) for \(0\le k<r\). Adding unused vertices as singleton parts does not change the sum, so assume \(\sum_i n_i=n\).

For fixed \(S\), the function \(f(x)+f(S-x)\) is strictly nondecreasing for \(x\ge \max\{S/2,r-1\}\), because
\[
[f(x+1)+f(S-x-1)]-[f(x)+f(S-x)]
=\binom{x}{r-1}-\binom{S-x-1}{r-1}>0.
\]
Consequently, if two nontrivial parts have sizes \(p\ge q\) and \(p<M\), replacing them by \(p+q\) when \(p+q\le M\), or by \(M\) and \(p+q-M\) when \(p+q>M\), increases the sum. A remainder between \(1\) and \(r-1\) is represented by singleton parts. Repeating this operation produces \(\alpha\) parts of size \(M\) and the remainder \(\beta\). The inequality follows.

Whenever this operation is applied to two nontrivial parts and the larger part is smaller than \(M\), the displayed difference is positive. Hence equality is possible only for the stated extremal configurations.
\end{proof}

\begin{lemma}\label{bounded-degree-clique-equality}
Let $M\ge r\ge3$, and let $H$ be a graph on $m$ vertices with $\Delta(H)\le M-1$ such that every edge of $H$ is contained in an $r$-clique. Write
\[
m=aM+b,
\qquad a\ge0,
\qquad 0\le b\le M-1.
\]
Then
\[
|\mathcal C_r(H)|
\le a\binom Mr+\binom br.
\]
Moreover, equality holds if and only if the nontrivial components of $H$ are
\[
aK_M\dunion K_b
\]
when $b\ge r$, and are exactly $a$ copies of $K_M$ when $b<r$; in the latter case the remaining $b$ vertices are isolated.
\end{lemma}

\begin{proof}
The inequality is proved by Chao and Dong~\cite{Chao2022}. We prove the equality statement from the argument of Chao and Dong~\cite{Chao2022}.

For an integer $x\ge0$, write $x=qM+\gamma$ with $0\le \gamma<M$ and set
\[
F(x)=q\binom Mr+\binom{\gamma}{r}.
\]
The same discrete convexity used in Lemma~\ref{opt} gives, for every $1\le k\le M$,
\begin{equation}\label{bounded-degree-convex-step}
F(m-k)+\binom kr\le F(m).
\end{equation}
Indeed, if $k\le b$, this reduces to
\[
\binom{b-k}{r}+\binom kr\le\binom br,
\]
while if $k>b$, it reduces to
\[
\binom kr+\binom{M+b-k}{r}
\le
\binom Mr+\binom br.
\]

Suppose now that equality holds in the stated clique bound. For $v\in V(H)$, let $\mathcal T_v$ be the family of $r$-cliques meeting the closed neighborhood $N_H[v]$. Chao and Dong~\cite{Chao2022} proved the local double-counting inequality (Lemma 2 in \cite{Chao2022})
\begin{equation}\label{Chao-Dong-local-sum}
\sum_{v\in V(H)}|\mathcal T_v|
\le
\sum_{v\in V(H)}\binom{d_H(v)+1}{r}.
\end{equation}
If~\eqref{Chao-Dong-local-sum} were strict, then some vertex $v$ would satisfy
\[
|\mathcal T_v|<\binom{d_H(v)+1}{r}.
\]
Putting $k=d_H(v)+1\le M$ and applying the bounded-degree clique theorem to $H-N_H[v]$, we would obtain, by~\eqref{bounded-degree-convex-step},
\[
\begin{aligned}
|\mathcal C_r(H)|
&=|\mathcal C_r(H-N_H[v])|+|\mathcal T_v|\\
&<F(m-k)+\binom kr\\
&\le F(m),
\end{aligned}
\]
a contradiction. Hence equality holds in~\eqref{Chao-Dong-local-sum}.

The proof of Lemma 2 in ~\eqref{Chao-Dong-local-sum}  implies that equality in~\eqref{Chao-Dong-local-sum} forces
\begin{equation}\label{bounded-degree-local-equality}
c_r(v)=\binom{d_H(v)}{r-1}
\qquad\text{for every }v\in V(H).
\end{equation}

Let $v$ be nonisolated. Since every edge of $H$ lies in an $r$-clique, the vertex $v$ lies in an $r$-clique and hence $d_H(v)\ge r-1$. Equality~\eqref{bounded-degree-local-equality} says that every $(r-1)$-subset of $N_H(v)$ forms a clique. Since $r\ge3$, this implies that $N_H(v)$ is complete. Consequently every nontrivial component of $H$ is complete. The maximum-degree condition shows that every such component has order at most $M$, and the $r$-clique-supported assumption shows that every nontrivial component has order at least $r$.

If $a=0$, the conclusion is immediate. Suppose $a\ge1$. If $M=r$, then every nontrivial component has order exactly $M$. Since $b<M=r$, the right-hand side of the clique bound equals $a$, while each nontrivial component contributes exactly one $r$-clique. Thus equality forces exactly $a$ components of order $M$, and the remaining $b$ vertices are isolated. We may therefore assume $M>r$. Applying Lemma~\ref{opt} to the orders of the complete nontrivial components, together with singleton parts for the isolated vertices, its equality statement yields exactly $a$ components of order $M$ and, when $b\ge r$, one further component of order $b$; when $b<r$, the remaining $b$ vertices are isolated.
\end{proof}

\begin{proof}[Proof of Theorems~\ref{leaf-bunch-counting-theorem} and~\ref{leaf-bunch-spectral-theorem}]
Let $G$ be an $n$-vertex $T_t$-free graph and put $d=t-r$. Let $v$ be a vertex contained in an $r$-clique. If $d_G(v)\ge t-1$, choose $d$ leaves with a common parent $p$ in $T_t$, embed the remaining $r$ vertices of $T_t$ into an $r$-clique containing $v$ with $p$ mapped to $v$, and map the deleted leaves to $d$ distinct neighbors of $v$ outside the clique. This produces a copy of $T_t$, a contradiction. Hence every vertex lying in an $r$-clique has degree at most $t-2$.

For the clique-counting statement, write $n=\alpha(t-1)+\beta$ as in Theorem~\ref{leaf-bunch-counting-theorem}, and let $G_r$ be the spanning subgraph of $G$ whose edges are precisely those contained in at least one $r$-clique. Then $\mathcal C_r(G_r)=\mathcal C_r(G)$ and $\Delta(G_r)\le t-2$. The Gan--Loh--Sudakov conjecture, proved by Chase~\cite{chase2020} and also by Chao and Dong~\cite{Chao2022}, gives
\[
|\mathcal C_r(G)|=|\mathcal C_r(G_r)|
\le \alpha\binom{t-1}{r}+\binom{\beta}{r},
\]
which gives the upper bound in~\eqref{leaf-bunch-counting}. Since $\alpha K_{t-1}\dunion K_\beta$ is $T_t$-free and has exactly
\[
\alpha\binom{t-1}{r}+\binom{\beta}{r}
\]
copies of $K_r$, equality in~\eqref{leaf-bunch-counting} follows. It remains to characterize the extremal graphs. Suppose that $G$ attains equality in~\eqref{leaf-bunch-counting}. By construction, every edge of $G_r$ lies in an $r$-clique. Applying Lemma~\ref{bounded-degree-clique-equality} with $M=t-1$ shows that, if $\beta\ge r$, then the nontrivial components of $G_r$ are $\alpha$ copies of $K_{t-1}$ together with one copy of $K_\beta$, whereas if $\beta<r$, the nontrivial components of $G_r$ are exactly $\alpha$ copies of $K_{t-1}$ and the remaining $\beta$ vertices are isolated in $G_r$. By Lemma~\ref{isolation}, every such $K_{t-1}$ is a component of the original graph $G$. Hence, when $\beta\ge r$, no further edge can be added and $G=\alpha K_{t-1}\dunion K_\beta$; when $\beta<r$, the remaining $\beta$ vertices may induce an arbitrary graph. Conversely, each graph of the stated form is $T_t$-free and attains equality. This proves Theorem~\ref{leaf-bunch-counting-theorem}.

For the clique-spectral statement, every vertex $v$ contained in an $r$-clique satisfies
\[
c_r(v)\le\binom{t-2}{r-1}
\]
since $d_G(v)\le t-2$.
Thus~\eqref{local-clique-degree} gives~\eqref{leaf-bunch-criterion-bound}. If equality holds, then some vertex $v$ satisfies $c_r(v)=\binom{t-2}{r-1}$. Consequently $d_G(v)=t-2$ and every $(r-1)$-subset of $N_G(v)$ is a clique. Since $r\ge3$, the neighborhood $N_G(v)$ is complete, so $N_G[v]=K_{t-1}$; Lemma~\ref{isolation} makes it a component. The converse is immediate. This proves Theorem~\ref{leaf-bunch-spectral-theorem}.
\end{proof}

\begin{proof}[Proof of Theorem~\ref{T5}]
Put $r=t-d$ and $\lambda=\lambda(T_t)$, and let $G$ be an extremal $n$-vertex $T_t$-free graph.

If $\lambda\ge d$, then $\lambda(T_t)\ge t-r$, so Theorem~\ref{leaf-bunch-counting-theorem} gives exactly the asserted generalized Tur\'{a}n number and the stated characterization of the extremal graphs.

Assume now that $\lambda\le d-1$. Let $G_r$ be obtained from $G$ by deleting every edge contained in no $r$-clique. This operation preserves all copies of $K_r$. Since $t\ge d^2-d+3$, Lemma~\ref{component-bound} shows that every nontrivial component of $G_r$ has order at most $t-1$. Let their orders be $n_1,\ldots,n_k$. Each $n_i\ge r$, and therefore
\[
\begin{aligned}
|\mathcal C_r(G)|
&=|\mathcal C_r(G_r)|\\
&\le\sum_{i=1}^k\binom{n_i}{r}\\
&\le\alpha\binom{t-1}{r}+\binom{\beta}{r},
\end{aligned}
\]
where the last inequality follows from Lemma~\ref{opt}. The graph $\alpha K_{t-1}\dunion K_\beta$ is $T_t$-free and attains this bound.

Suppose equality holds. Then every nontrivial component of $G_r$ must be complete, and the equality statement in Lemma~\ref{opt} gives the component orders stated in the theorem. Lemma~\ref{isolation} shows that every copy of $K_{t-1}$ is a component of the original graph $G$. If $\beta<r$, the remaining $\beta$ vertices contain no counted clique and may induce an arbitrary graph. If $\beta\ge r$, they must induce $K_\beta$. This proves the theorem.
\end{proof}



\subsection{Clique-spectral proofs}\label{spectral-proofs}

We first prove two spectral lemmas used in the proof of Theorem~\ref{spectral-T5}.

\begin{lemma}\label{complete-spectral-bound}
Let $H$ be a graph on $m\ge r$ vertices. Then
\[
\rho_r(H)\le \rho_r(K_m)=\binom{m-1}{r-1}.
\]
Moreover, equality holds if and only if $H=K_m$.
\end{lemma}

\begin{proof}
Let $x=(x_1,\ldots,x_m)\in\mathbb R_+^m$ satisfy $\sum_{i=1}^m x_i^r=1$, and let $e_r(x)=\sum_{1\le i_1<i_2<\ldots <i_r\le m}x_{i_1}\cdots x_{i_r}$
denote the $r$-th elementary symmetric function of $x_1,\ldots,x_m$. By~\eqref{spectral-variational},
\[
r\sum_{\{i_1,\ldots,i_r\}\in\mathcal C_r(H)}x_{i_1}\cdots x_{i_r}
\le r e_r(x).
\]
Maclaurin's inequality and H\"older's inequality give
\[
\left(\frac{e_r(x)}{\binom mr}\right)^{1/r}
\le \frac{x_1+\cdots+x_m}{m}
\le m^{-1/r}.
\]
Consequently,
\[
r e_r(x)\le \frac{r}{m}\binom mr=\binom{m-1}{r-1}.
\]
Taking the maximum over all such $x$ proves the desired upper bound. Equality is attained for $K_m$ by taking $x_1=\cdots=x_m=m^{-1/r}$, so
\[
\rho_r(K_m)=\binom{m-1}{r-1}.
\]
If equality holds for $H$, then equality must hold in H\"older's inequality, and hence all $x_i$ are equal and positive. Equality in the first displayed inequality then implies that every $r$-subset of $V(H)$ is an $r$-clique. Since $m\ge r\ge2$, this forces $H=K_m$.
\end{proof}

\begin{lemma}\label{spectral-component-max}
Let $r\ge2$, and let $H$ be a graph whose nontrivial components containing an $r$-clique are $C_1,\ldots,C_k$. If $H$ contains at least one $r$-clique, then
\[
\rho_r(H)=\max_{i\in[k]}\rho_r(C_i).
\]
\end{lemma}

\begin{proof}
Let $x\ge0$ be admissible in~\eqref{spectral-variational}, and put
\[
a_i=\sum_{v\in V(C_i)}x_v^r.
\]
After normalizing the restriction of $x$ to $C_i$, the contribution of $C_i$ to~\eqref{spectral-variational} is at most $a_i\rho_r(C_i)$. Vertices outside these components contribute nothing. Hence the total is at most
\[
\sum_i a_i\rho_r(C_i)\le \max_i\rho_r(C_i).
\]
The reverse inequality follows by taking a maximizing vector supported on one component.
\end{proof}

\begin{proof}[Proof of Theorem~\ref{spectral-T5}]
Put $r=t-d$ and $\lambda=\lambda(T_t)$. Notice that
\[
t\ge d^2-d+3\ge3d-1,
\]
and hence $r\ge3$.

If $\lambda\ge d$, Theorem~\ref{leaf-bunch-spectral-theorem} gives
\[
\rho_r(G)\le\binom{t-2}{r-1}.
\]
Its equality statement gives exactly the required characterization.

Assume now that $\lambda\le d-1$. Let $G_r$ be the spanning subgraph of $G$ whose edges are precisely those contained in at least one $r$-clique. Deleting edges contained in no $r$-clique does not change the $r$-clique tensor, so
\[
\rho_r(G)=\rho_r(G_r).
\]
If $G_r$ has no $r$-clique, then $\rho_r(G)=0$ and the result is immediate. Otherwise let $C_1,\ldots,C_k$ be the nontrivial components of $G_r$.  By Lemma~\ref{component-bound}, every nontrivial component of $G_r$ has order at most $t-1$. By Lemma~\ref{spectral-component-max} and Lemma~\ref{complete-spectral-bound},
\[
\rho_r(G)
=\max_i\rho_r(C_i)
\le\max_i\binom{|V(C_i)|-1}{r-1}
\le\binom{t-2}{r-1}.
\]
If equality holds, some $C_i$ has order $t-1$ and equality holds in Lemma~\ref{complete-spectral-bound}; hence $C_i=K_{t-1}$. Lemma~\ref{isolation} shows that this copy is a component of $G$.
 This completes the proof.
\end{proof}


\section{Concluding remarks and open problems}\label{concluding-remarks}

 Gan, Loh and Sudakov\cite{gan2015maximizing} showed that if Conjecture~\ref{C1} holds for  \(r\) and the star, then it holds for $r+1$ and the star. We conjecture that this holds for any tree.




\begin{conjecture}\label{C2}
If Conjecture~\ref{C1} holds for a given \(r\), then it also holds for \(r+1\).
\end{conjecture}




There is a more global way to place Theorem~\ref{CDT-spectral-Erdos-Sos} and our large-clique theorem (Theorem~\ref{spectral-T5}) in one framework. Let $\mathcal T_t$ denote the family of all trees on $t$ vertices.  For $2\le r\le t-1$, define
\begin{equation}\label{universal-spectral-extremal}
\Lambda_r(n,t)
=
\max\left\{
\rho_r(G): |V(G)|=n,\ \text{$G$ does not contain every member of $\mathcal T_t$}
\right\}.
\end{equation}
Equivalently, a graph is admissible in~\eqref{universal-spectral-extremal} if it misses at least one $t$-vertex tree. Thus $\Lambda_r(n,t)$ is the natural $r$-clique spectral analogue of the universal adjacency-spectral extremal problem considered by Cioab\u{a}, Desai and Tait.

For $r=2$, Theorem~\ref{CDT-spectral-Erdos-Sos} is equivalent to the following exact values, together with uniqueness of the extremal graph: for every fixed $k\ge2$ and all sufficiently large $n$,
\begin{align}
\Lambda_2(n,2k+2)&=\rho(S_{n,k}),\label{Lambda-CDT-even}\\
\Lambda_2(n,2k+3)&=\rho(S_{n,k}^{+}).\label{Lambda-CDT-odd}
\end{align}
At the other end, Theorems~\ref{spectral-T1} and~\ref{spectral-T5} give the solution for large clique orders. In particular, if $d=t-r\ge2$, $t\ge d^2-d+3$, then
\begin{equation}\label{Lambda-large-clique}
\Lambda_r(n,t)=\rho_r(K_{t-1})=\binom{t-2}{r-1}
\end{equation}
for every $n\ge t-1$; the same formula also holds for $r=t-1$ by Theorem~\ref{spectral-T1}. Moreover, Theorem~\ref{spectral-T5} characterizes all extremal admissible graphs in the range $d\ge2$.

The two constructions naturally occupy different ranges of clique orders. If $2\le r\le k+1$, fix $r-1$ vertices in the $K_k$-part of $S_{n,k}$. Together with each vertex of the independent set they form an $r$-clique. Substituting into~\eqref{spectral-variational} a vector which has value $r^{-1/r}$ on these $r-1$ core vertices and value $(r(n-k))^{-1/r}$ on the independent set gives
\begin{equation}\label{split-lower-bound}
\rho_r(S_{n,k})\ge (n-k)^{(r-1)/r}.
\end{equation}
The same lower bound holds for $S_{n,k}^{+}$. Hence, for fixed $k$ and $r\le k+1$, the clique-spectral radius of the split construction tends to infinity with $n$, whereas
\[
\rho_r(K_{t-1})=\binom{t-2}{r-1}
\]
is independent of $n$. On the other hand,
\[
\omega(S_{n,k})=k+1,
\qquad
\omega(S_{n,k}^{+})=k+2.
\]
Thus $S_{n,k}$ has no $r$-cliques for $r\ge k+2$. Moreover, $S_{n,k}^{+}$ has exactly one $(k+2)$-clique, so
\[
\rho_{k+2}(S_{n,k}^{+})=1<\binom{2k+1}{k+1}=\rho_{k+2}(K_{2k+2}),
\]
and it has no $r$-cliques for $r\ge k+3$. These observations lead to the following conjecture.

\begin{conjecture}[Clique-spectral extension of the Cioab\u{a}--Desai--Tait theorem]\label{universal-clique-spectral-conjecture}
Fix $k\ge2$ and an integer $r\ge2$. For all sufficiently large $n$, the following statements hold.
\begin{enumerate}[(i)]
\item For trees of even order $t=2k+2$,
\[
\Lambda_r(n,2k+2)=
\begin{cases}
\rho_r(S_{n,k}),&2\le r\le k+1,\\[1mm]
\rho_r(K_{2k+1})=\displaystyle\binom{2k}{r-1},&k+2\le r\le2k+1.
\end{cases}
\]
In the first range, $S_{n,k}$ is the unique extremal graph. In the second range, every extremal admissible graph contains $K_{2k+1}$ as a component.

\item For trees of odd order $t=2k+3$,
\[
\Lambda_r(n,2k+3)=
\begin{cases}
\rho_r(S_{n,k}^{+}),&2\le r\le k+1,\\[1mm]
\rho_r(K_{2k+2})=\displaystyle\binom{2k+1}{r-1},&k+2\le r\le2k+2.
\end{cases}
\]
In the first range, $S_{n,k}^{+}$ is the unique extremal graph. In the second range, every extremal admissible graph contains $K_{2k+2}$ as a component.
\end{enumerate}
\end{conjecture}





\section*{Acknowledgments}
We thank Daniel Gerbner for valuable communications. This research is supported in part by the National Natural Science Foundation of China (Nos.~12571363 and 12371327) and the Natural Science Foundation of Hunan Province (Grant No.~2025JJ30003).

\section*{Declaration on the use of generative AI}
During the preparation of this manuscript, the authors used GPT 5.5 to assist with English-language editing. The authors take full responsibility for the mathematical content.


\begin{thebibliography}{99}

\bibitem{ajtai1994} M. Ajtai and J. Koml\'{o}s and E. Szemeredi, On a conjecture of Loebl, (1994).

\bibitem{alexander2012} J. Alexander and J. Cutler and T. Mink, Independent sets in graphs with given minimum degree, Electr. J. Comb. 19 (2012), P37--P37.

\bibitem{alexander2016} J. Alexander and T. Mink, A new method for enumerating independent sets of a fixed size in general
 graphs, J. Graph Theory 81 (2016), 57--72.

\bibitem{alon2016}
N. Alon and C. Shikhelman,
Many \(T\) copies in \(H\)-free graphs,
\emph{J. Combin. Theory Ser. B} 121 (2016), 146--172.

\bibitem{besomi2019}
G. Besomi, M. Pavez-Sign\'{e}, and M. Stein,
Degree conditions for embedding trees,
\emph{SIAM J. Discrete Math.} 33 (2019), 1521--1555.

\bibitem{Brennan2016}
M. Brennan,
Ramsey numbers of trees versus odd cycles,
\emph{Electron. J. Combin.} 23(3) (2016), Paper No.~P3.2.

\bibitem{Brennan2017}
M. Brennan,
Ramsey numbers of trees and unicyclic graphs versus fans,
\emph{Discrete Math.} 340(5) (2017), 969--983.

\bibitem{Brandt1996}
S. Brandt and E. Dobson,
The Erd\H{o}s--S\'{o}s conjecture for graphs of girth 5,
\emph{Discrete Math.} 150 (1996), 411--414.

\bibitem{chakraborti2024exact}
D. Chakraborti and D. Chen,
Exact results on generalized Erd\H{o}s--Gallai problems,
\emph{European J. Combin.} 120 (2024), 103955.

\bibitem{Chao2022}
T.-W. Chao and Z. Dong,
A simple proof of the Gan--Loh--Sudakov conjecture,
\emph{Electron. J. Combin.} 29(3) (2022), P3.59.

\bibitem{chase2020}
Z. Chase,
The maximum number of triangles in a graph of given maximum degree,
\emph{Adv. Comb.} (2020), Paper No.~10.

\bibitem{CioabaDesaiTait2023}
S. Cioab\u{a}, D. N. Desai, and M. Tait,
A spectral Erd\H{o}s--S\'{o}s theorem,
\emph{SIAM J. Discrete Math.} 37 (2023), 2228--2239.

\bibitem{cutler2014maximum}
J. Cutler and A. J. Radcliffe,
The maximum number of complete subgraphs in a graph with given maximum degree,
\emph{J. Combin. Theory Ser. B} 104 (2014), 60--71.

\bibitem{engbers2014counting} J. Engbers and D. Galvin, Counting independent sets of a fixed size in graphs with a given minimum degree, J. Graph Theory. 76(2014), 149--168.

\bibitem{erdos1962}
P. Erd\H{o}s,
Extremal problems in graph theory,
in \emph{Theory of Graphs and Its Applications}, Publishing House of the Czechoslovak Academy of Sciences, Prague, 1964, 29--36.

\bibitem{1959Erdos}
P. Erd\H{o}s and T. Gallai,
On maximal paths and circuits of graphs,
\emph{Acta Math. Acad. Sci. Hungar.} 10 (1959), 337--356.

\bibitem{ES66} P. Erd\H{o}s, M. Simonovits, A limit theorem in graph theory, Stud. Sci. Math. Hungar. 1 (1966) 51--57.

\bibitem{ES46} P. Erd\H{o}s, A. H. Stone, On the structure of linear graphs, Bull. Amer. Math. Soc. 52 (1946) 1087--1091.

\bibitem{Fan2013} G. Fan, The Erd{\H{o}}s--S{\'o}s conjecture for spiders of large size, Discrete Math. 313 (2013), 2513--2517.

\bibitem{Fan2007} G. Fan and L. Sun, The Erd{\H{o}}s--S{\'o}s conjecture for spiders, Discrete Math. 307 (2007), 3055--3062.

\bibitem{Fan2019}
G. Fan, Y. Hong, and Q. Liu,
The Erd{\H{o}}s--S{\'o}s conjecture for spiders,
\emph{arXiv preprint} arXiv:1804.06567v2 (2019).

\bibitem{GMP2020}
D. Gerbner, A. Methuku, and C. Palmer,
General lemmas for Berge--Tur\'{a}n hypergraph problems,
\emph{European J. Combin.} 86 (2020), 103082.

\bibitem{GP2026}
D. Gerbner and C. Palmer,
Survey of generalized Tur\'{a}n problems---counting subgraphs,
\emph{Electron. J. Combin.}, Dynamic Surveys (2026), DS27.

\bibitem{gerbner2023doublestars}
D. Gerbner,
Generalized Tur\'{a}n problems for double stars,
\emph{Discrete Math.} 346(7) (2023), 113395.

\bibitem{goerlich2016}
A. Goerlich and A. \.{Z}ak,
On the Erd\H{o}s--S\'{o}s conjecture for trees of large size,
\emph{Electron. J. Combin.} 23 (2016), P1.52.

\bibitem{gan2015maximizing}
W. Gan, P.-S. Loh, and B. Sudakov,
Maximizing the number of independent sets of a fixed size,
\emph{Combin. Probab. Comput.} 24 (2015), 521--527.

\bibitem{Law2013} H. Law and C. McDiarmid, On independent sets in graphs with given minimum degree, Combin. Probab. Comput. 22(2013), 874--884.

\bibitem{letzter2019}
S. Letzter,
Many \(H\)-copies in graphs with a forbidden tree,
\emph{SIAM J. Discrete Math.} 33(4) (2019), 2360--2368.

\bibitem{luo2018maximum}
R. Luo,
The maximum number of cliques in graphs without long cycles,
\emph{J. Combin. Theory Ser. B} 128 (2018), 219--226.

\bibitem{rozhon2019}
V. Rozho\v{n},
A local approach to the Erd\H{o}s--S\'{o}s conjecture,
\emph{SIAM J. Discrete Math.} 33 (2019), 643--664.

\bibitem{sacle1997}
J. Sacl\'{e} and M. W\'{o}zniak,
A note on the Erd\H{o}s--S\'{o}s conjecture for graphs without \(C_4\),
\emph{J. Combin. Theory Ser. B} 70 (1997), 229--234.

\bibitem{LiuBu2023}
C. Liu and C. Bu,
On a generalization of the spectral Mantel's theorem,
\emph{J. Combin. Optim.} 46 (2023), Paper No.~14.

\bibitem{Nikiforov2010paths}
V. Nikiforov,
The spectral radius of graphs without paths and cycles of specified length,
\emph{Linear Algebra Appl.} 432 (2010), 2243--2256.

\bibitem{Qi2013}
L. Qi,
Symmetric nonnegative tensors and copositive tensors,
\emph{Linear Algebra Appl.} 439 (2013), 228--238.

\bibitem{t1941}
P. Tur\'{a}n,
On an extremal problem in graph theory,
\emph{Mat. Fiz. Lapok} 48 (1941), 436--452.

\bibitem{ZhouYuan2026}
J. Zhou and X. Yuan,
Counting large cliques in graphs with a forbidden tree,
\emph{arXiv preprint} arXiv:2607.23960 (2026).

\bibitem{ZP1} X. Zhao and Y. Peng, The maximum number of triangles in graphs without cycles of length $0\mod5$, \emph{Journal of Graph Theory}, https://doi.org/10.1002/jgt.70111.

\bibitem{zykov1949}
A. A. Zykov,
On some properties of linear complexes,
\emph{Mat. Sbornik} 24(66) (1949), 163--188.

\end{thebibliography}
\end{document}
